\chapter*{Résumé}
\addcontentsline{toc}{chapter}{Résumé}

Ce rapport présente le travail réalisé dans le cadre du Projet de Fin d'Année (PFA) de 2\textsuperscript{e} année du cycle ingénieur en Génie Logiciel à l'ISIMS de Sfax. Le stage, effectué à distance au sein de \textbf{SINAPS TEC} (Solutions Innovations For Networks, Applications \& Portable Systems) du 22~juin au 6~septembre~2026, a consisté à concevoir et développer \textbf{Op Support IA}, une application web de support client assisté par un agent d'intelligence artificielle.

Le système permet à un client d'ouvrir une conversation, d'obtenir une réponse automatique fondée sur une base de connaissances (approche RAG) et, si besoin, de basculer vers un agent humain en temps réel. Des suggestions d'action rapide accompagnent les réponses IA, tandis que les messages, horodatages, avatars et changements de statut sont persistés. Un espace agent gère la file d'attente des demandes escaladées ; un espace administrateur valide les comptes, consulte l'historique et affiche des indicateurs (volume, résolution IA ou humaine, satisfaction, délai de réponse).

La solution comprend un frontend web \textbf{Next.js} (React, TypeScript), un client mobile \textbf{Expo / React Native} et un backend \textbf{Node.js / Express}, avec persistance \textbf{MongoDB}, authentification JWT et Google, génération de réponses via \textbf{Gemini} et un retriever local TF-IDF / similarité cosinus. Des tests automatisés (Jest, Supertest) couvrent le backend ; le client mobile possède en outre une vérification de typage TypeScript, sans suite de tests automatisés ni recette formelle documentée dans le dépôt.

\vspace{8pt}
\noindent\textbf{Mots-clés :} support client, agent IA, RAG, Next.js, Express, MongoDB, Socket.io, JWT, PFA.

\vspace{1.4cm}
\chapter*{Résumé en anglais}
\addcontentsline{toc}{chapter}{Résumé en anglais}

Ce rapport décrit un projet de stage de deuxième année du cycle ingénieur en Génie Logiciel (PFA) réalisé à l'ISIMS de Sfax. Le stage à distance au sein de \textbf{SINAPS TEC} s'est déroulé du 22~juin au 6~septembre~2026. Le système obtenu, \textbf{Op Support IA}, comprend une application web et un client mobile de support client assisté par IA.

Le client peut ouvrir une conversation, recevoir une réponse automatique fondée sur une base de connaissances (RAG) et demander l'intervention d'un agent humain lorsque cela est nécessaire. Les agents traitent la file des demandes escaladées ; les administrateurs valident les comptes, consultent l'historique et examinent des indicateurs opérationnels orientés SLA.

L'implémentation utilise un frontend \textbf{Next.js}, un client mobile \textbf{Expo / React Native} et un backend \textbf{Node.js / Express}, avec \textbf{MongoDB}, une authentification JWT et Google, \textbf{Gemini} associé à un retriever local TF-IDF / similarité cosinus, ainsi que des tests automatisés avec Jest et Supertest.

\vspace{8pt}
\noindent\textbf{Mots-clés :} support client, agent IA, RAG, Next.js, React Native, Express, MongoDB, Socket.io, JWT.
